\documentclass[12pt]{amsart}
%prepared in AMSLaTeX, under LaTeX2e
\addtolength{\oddsidemargin}{-0.6in} 
\addtolength{\evensidemargin}{-0.6in}
\addtolength{\topmargin}{-.55in}
\addtolength{\textwidth}{1.3in}
\addtolength{\textheight}{.9in}

\renewcommand{\baselinestretch}{1.05}

\usepackage{verbatim} % for "comment" environment

\usepackage{palatino}

\newtheorem*{thm}{Theorem}
\newtheorem*{defn}{Definition}
\newtheorem*{example}{Example}
\newtheorem*{problem}{Problem}
\newtheorem*{remark}{Remark}

\usepackage{fancyvrb,xspace,dsfont}

\usepackage[final]{graphicx}

% macros
\usepackage{amssymb}

\usepackage{hyperref}
\hypersetup{pdfauthor={Ed Bueler},
            pdfcreator={pdflatex},
            colorlinks=true,
            citecolor=blue,
            linkcolor=red,
            urlcolor=blue,
            }

\newcommand{\bn}{\mathbf{n}}
\newcommand{\br}{\mathbf{r}}
\newcommand{\bt}{\mathbf{t}}
\newcommand{\bu}{\mathbf{u}}
\newcommand{\bv}{\mathbf{v}}
\newcommand{\bx}{\mathbf{x}}
\newcommand{\by}{\mathbf{y}}

\newcommand{\cB}{\mathcal{B}}
\newcommand{\cD}{\mathcal{D}}
\newcommand{\cF}{\mathcal{F}}
\newcommand{\cH}{\mathcal{H}}
\newcommand{\cL}{\mathcal{L}}
\newcommand{\cV}{\mathcal{V}}
\newcommand{\cW}{\mathcal{W}}

\newcommand{\CC}{\mathbb{C}}
\newcommand{\NN}{\mathbb{N}}
\newcommand{\RR}{\mathbb{R}}
\newcommand{\ZZ}{\mathbb{Z}}

\renewcommand{\Im}{\mathrm{Im}}
\renewcommand{\Re}{\mathrm{Re}}

\newcommand{\eps}{\epsilon}
\newcommand{\grad}{\nabla}
\newcommand{\lam}{\lambda}
\newcommand{\lap}{\triangle}

\newcommand{\ip}[2]{\ensuremath{\left<#1,#2\right>}}

\newcommand{\image}{\operatorname{im}}
\newcommand{\onull}{\operatorname{null}}
\newcommand{\rank}{\operatorname{rank}}
\newcommand{\range}{\operatorname{range}}
\newcommand{\trace}{\operatorname{tr}}
\newcommand{\Span}{\operatorname{span}}

\newcommand{\prob}[1]{\bigskip\noindent\textbf{#1.}\quad }

\newcommand{\epart}[1]{\medskip\noindent\textbf{(#1)}\quad }
\newcommand{\ppart}[1]{\,\textbf{(#1)}\quad }

\newcommand{\ds}{\displaystyle}

\newcommand{\snew}{s^{\text{new}}}
\newcommand{\rhoi}{\rho_{\text{i}}}



\begin{document}
\scriptsize \hfill \emph{June 2026 \,(Bueler)}
\normalsize\medskip

\Large\centerline{\textbf{Solutions to Paper Exercises}}

\normalsize
\medskip
\begin{quote}
\emph{Acronyms: SKE = surface kinematical equation, SIA = shallow ice approximation, CFL = Courant, Friedrichs, Lewy (stability criterion).}
\end{quote}

\prob{1}  One way to think about normal vectors is to first construct the general form of a tangent vector.  In 2D we have a curve $z=s(x)$ and in 3D a surface $z=s(x,y)$.  In the first case, if we change the $x$ value by any number $\Delta x$ then by the linearization (differential) of $s$ we have $\Delta s= \frac{ds}{dx} \Delta x$ as the change in the $z$ direction.  Thus any tangent vector is
	$$\bt_s = \left(\Delta x, \frac{ds}{dx} \Delta x\right)$$
for some value of $\Delta x$.  On the other hand, for the given vector, ignoring time, we have $\bn_s = (-\frac{ds}{dx},1)$, thus
	$$\bn_s \cdot \bt_s = -\Delta x \frac{ds}{dx} + \frac{ds}{dx} \Delta x = 0$$
so $\bn_s$ is orthogonal ($=$normal) to any tangent vector.

In 3D, for the surface $z=s(x,y)$, the equivalent formulas are
\begin{align*}
\bt_s &= \left(\Delta x, \Delta y, \frac{\partial s}{\partial x} \Delta x + \frac{\partial s}{\partial y} \Delta y\right) \quad \text{ and } \quad
\bn_s = \left(-\frac{\partial s}{\partial x},-\frac{\partial s}{\partial y},1\right)
\end{align*}
Note that in $\bt_s$ one can choose any values $\Delta x,\Delta y$; there is a tangent \emph{plane}.  It is easy to check that $\bn_s \cdot \bt_s = 0$.


\prob{2}  The forward-time centered-space (FTCS) scheme in the slides is
	$$\snew_j = s_j - \Delta t\, u_j \frac{s_{j+1}-s_{j-1}}{2\Delta x} + \Delta t (a_j + w_j).$$
Assume $a=0$ and $w=0$, and write in the suggested form:
	$$\snew_j = \underbrace{\left(\frac{u_j \Delta t}{2\Delta x}\right)}_{c_{-}} s_{j-1} + \underbrace{\phantom{\frac{1}{1}} 1 \phantom{\frac{1}{1}}}_{c_0}\, s_j + \underbrace{\left(-\frac{u_j \Delta t}{2\Delta x}\right)}_{c_{+}} s_{j+1}.$$
Thus, as claimed,
	$$c_-+c_0+c_+ = \frac{u_j \Delta t}{2\Delta x} + 1 - \frac{u_j \Delta t}{2\Delta x} = 1.$$

However, depending on the sign of the velocity $u_j$, one of the coefficients $c_-,c_+$ is negative, so the scheme does not compute an actual average of the values $s_{j-1},s_j,s_{j+1}$.  For example, if $u_j >0$ then $c_+<0$.  For $A=0$, $B=0$, and $C=1$ this ``average'' would be then be the negative value
	$$c_- A + c_0 B + c_+ C = c_+ < 0,$$
but none of $A,B,C$ are negative.  If $u_j<0$ then $c_-<0$ and a similar example could be constructed.  On the other hand, an actual averaging process would have the property that the average of $A,B,C$ was somewhere between the smallest and the largest of these numbers, which is not true in this case.

\medskip \noindent \emph{Comment.}  Let us consider the significance of this non-averaging property of the FTCS scheme.  Suppose that the current gridded surface values $\{s_j\}_{j=0}^n$ have a ``wiggle'' or oscillation.  Then the new values $\{\snew_j\}_{j=0}^n$ can have oscillations which are a little bit bigger, because the new values can be outside of the range of the old values as shown above.  In this sense the scheme can be unstable, and in practice it is.  Careful consideration of this behavior uses \href{https://en.wikipedia.org/wiki/Von_Neumann_stability_analysis}{\emph{von Neumann stability analysis}}, a kind of Fourier analysis invented by the one who brought you the stored-program digital computer.  John von Neumann's idea was to substitute a pure gridded wave (\emph{mode}) into the scheme and measure the amount of growth/decay generated by one time step.  This analysis of the FTCS scheme for advection, as here, shows growth factors greater than one, whereas for the upwind scheme in the next exercise the growth factors are less than one.


\prob{3}  The upwind scheme for the SKE is in the slides.  Again assume $a_j=0$ and $w_j=0$.  For non-negative horizontal velocity $u_j\ge 0$ the scheme is:
    $$\snew_j = s_j - \Delta t\, u_j \frac{s_j-s_{j-1}}{\Delta x}$$
Collecting terms we get constants in the suggested form:
	$$\snew_j = \underbrace{\left(\frac{u_j \Delta t}{\Delta x}\right)}_{c_{-}} s_{j-1} + \underbrace{\left(1 - \frac{u_j \Delta t}{\Delta x}\right)}_{c_0} s_j$$
If $u_j\ge 0$ and $|u_j|\Delta t = u_j\Delta t\le \Delta x$ then both constants are nonnegative, $c_-\ge 0$ and $c_0\ge 0$.  (We are also assuming $\Delta t>0$ and $\Delta x > 0$.)  Furthermore, they add to one:
	$$c_{-} + c_0 = \frac{u_j \Delta t}{\Delta x} + 1 - \frac{u_j \Delta t}{\Delta x} = 1.$$
It follows that the new surface value $\snew_j$ is a true average of the current surface elevation pair $s_{j-1},s_j$.

A similar result holds for leftward flow.  For $u_j<0$, assuming CFL $|u_j|\Delta t = -u_j\Delta t \le \Delta x$,  $\snew_j = c_0 s_j + c_+ s_{j+1}$ is a strict average with $c_0\ge 0$, $c_+\ge 0$, and $c_0+c_+=1$.

This upwind method is stable \emph{if one uses time steps $\Delta t$ which satisfy CFL}.  This is because, unlike in exercise \textbf{2}, unstable modal growth is impossible.  If the current surface is any wave of any amplitude, the updated surface may still be wavy but the amplitude will be reduced.  Over time, all the bumps will be averaged-out.


\prob{4}  The requested derivative is a partial derivative because the integral depends on $t$.  Otherwise it is just a matter of applying the Leibniz rule directly:
{\small
\begin{align*}
\frac{\partial}{\partial x}\left(\int_{b(t,x)}^{s(t,x)} u(t,x,z)\,dz\right) &= \frac{\partial s}{\partial x}(t,x) u(t,x,s(t,x)) - \frac{\partial b}{\partial x}(x) u(t,x,b(x)) + \int_{b(t,x)}^{s(t,x)} \frac{\partial u}{\partial x}(t,x,z)\,dz
\end{align*}
}


\prob{5}  In what follows we will assume:
\begin{itemize}
\item the SKE holds,
\item the ice is incompressible,
\item the basal ice is non-sliding, and
\item the basal ice is non-penetrating.
\end{itemize}
The result of the derivation will be the mass continuity equation.

Recall that $H=s-b$.  We start the calculuation with the flux term in the mass continuity equation, expand the definition of $\bar u H$, then apply the Leibniz rule (\textbf{4} above), and then apply the basal condition on $u$:
\begin{align*}
\frac{\partial}{\partial x}(\bar u H) &= \frac{\partial}{\partial x}\left(\int_{b(t,x)}^{s(t,x)} u(t,x,z)\,dz\right) \\
   &= \frac{\partial s}{\partial x}(t,x) u(t,x,s(t,x)) - \frac{\partial b}{\partial x}(x) u(t,x,b(x)) + \int_{b(t,x)}^{s(t,x)} \frac{\partial u}{\partial x}(t,x,z)\,dz \\
   &= u(t,x,s(t,x))\,\frac{\partial s}{\partial x}(t,x) + \int_{b(t,x)}^{s(t,x)} \frac{\partial u}{\partial x}(t,x,z)\,dz.
\end{align*}
(The $u \partial s/\partial x$ term in the SKE has already appeared.)  Next we use incompressibility to rewrite the $x$ derivative of $u$ as the negative $z$ derivative of $w$:
  $$\frac{\partial}{\partial x}(\bar u H) = u(t,x,s(t,x))\,\frac{\partial s}{\partial x}(t,x) - \int_{b(t,x)}^{s(t,x)} \frac{\partial w}{\partial z}(t,x,z)\,dz$$
Now apply the fundamental theorem of calculus and the non-penetrating condition:
\begin{align*}
\frac{\partial}{\partial x}(\bar u H) &= u(t,x,s(t,x))\,\frac{\partial s}{\partial x}(t,x) - w(t,x,s(t,x)) + w(t,x,b(t,x)) \\
  &= u(t,x,s(t,x))\,\frac{\partial s}{\partial x}(t,x) - w(t,x,s(t,x))
\end{align*}
Thus we have turned the flux derivative into exactly the dynamic terms which appear in the SKE.  From now on we simplify appearance by hiding the $t,x$ dependence.

Because $\partial b/\partial t=0$ by assumption, we know $\partial s/\partial t = \partial H/\partial t$.  Also, we can write the SKE as $u \frac{\partial s}{\partial x} - w = -\frac{\partial s}{\partial t} + a$.  Thus
	$$\frac{\partial}{\partial x}(\bar u H) = u \frac{\partial s}{\partial x} - w = -\frac{\partial s}{\partial t} + a = -\frac{\partial H}{\partial t} + a$$
This is the requested mass continuity equation.


\prob{6}  Starting from the 5 boxed equations on slide 24, the formulas for $w(z)$, $p(z)$, and $\tau_{13}(z)$ on slide 25 are straightforward to get by integration, using the given boundary conditions.  Next, using the boundary condition $u(0)=u_0$ and the formula for $\tau_{13}(z)$ we have
    $$u(z) = u_0 + 2 A (\rho g \sin\alpha)^3 \int_0^z (H-z')^3\,dz'$$
Use the substitution $\zeta = H-z'$ and get the desired horizontal velocity formula:
\begin{align*}
u(z) &= u_0 - 2 A (\rho g \sin\alpha)^3 \int_H^{H-z} \zeta^3\,d\zeta \\
     &= u_0 - \frac{2}{4} A (\rho g \sin\alpha)^3  \left((H-z)^4 - H^4\right) \\
     &= u_0 + \frac{1}{2} A (\rho g \sin\alpha)^3  \left(H^4 - (H-z)^4\right)
\end{align*}

The claim is that the SIA uses this slab-on-a-slope formula, but with the local, i.e.~$x$-dependent, surface slope and thickness.  So, again using $\tan\alpha \approx \sin\alpha$, we should regard ``$\sin\alpha$'' as the local slope $\partial s/\partial x$.  On the other hand, we want to consider both positive and negative values for $\sin\alpha$, and we observe that the ice flows downhill whether it is flowing left or flowing right.  So if $m$ is the slope we want the 3rd power to behave properly:
	$$(-m)^3 \qquad \stackrel{\text{replace}}{\to} \qquad -|m|^2 m.$$

Thus we have the non-sliding, isothermal SIA formula on slide 28:
    $$u(t,x,z) = u_0 - \frac{1}{2} A (\rho g)^3 \left((s-b)^4 - (s-z)^4\right) \left|\frac{\partial s}{\partial x}\right|^2 \frac{\partial s}{\partial x}.$$
In this expression $s=s(t,x)$, $H=H(t,x)$, and $u_0(t,x)$.  However, the expression only makes sense at points which are inside the ice:
	$$b(x) < z < s(t,x)$$

\medskip
\noindent \emph{Comment.}  Before trusting this approximation of the solution of the power-law Stokes equations, observe that one does not actually get any basal sliding function $u_0(t,x)$ as data!  Rather, one needs to solve stress balances \emph{both} for deformation within the ice (e.g.~$\partial u/\partial z$) \emph{and} for the basal sliding velocity $u_0$.  A glaciological ``sliding law'' does not provide a known function $u_0$, but rather it relates $u_0$ to the basal stresses which one is also solving for.  In this context, however, there is no way for the SIA simplified model to balance the longitudinal (or ``membrane'') stresses which control the actual sliding of glaciers.  Therefore 21st-century practitioners recommend that the sliding be determined by solving a non-trivial stress balance.  One might use a hybrid balance, combining the SIA and the shallow shelf approximation for example, or a unified balance like the Blatter-Pattyn equations, or even the (full) Stokes model.


\prob{7}  From exercise \textbf{6}, for $n=3$ the surface value of the horizontal velocity in the non-sliding SIA is
\begin{align*}
u &= u(t,x,s(t,x)) = 0 - \frac{1}{2} A (\rho g)^3 \left((s-b)^4 - 0^4\right) \left|\frac{\partial s}{\partial x}\right|^2 \frac{\partial s}{\partial x} \\
  &= - \frac{1}{2} A (\rho g)^3 H^4 \left|\frac{\partial s}{\partial x}\right|^2 \frac{\partial s}{\partial x}
\end{align*}
The basic point, when one is thinking about how the non-sliding SIA model converts a given glacier shape into a surface velocity field, is that ice flows downhill in this model.  That is, the SIA model always makes the horizontal component point in the opposite direction from the surface gradient: $u \sim - \partial s/\partial x$.

The next point is that the magnitude $|u|$ scales with the 4th power of the ice thickness $H$.  For a given surface slope, relatively small increases in thickness correspond to substantial increases in surface speed.

On the other hand, the behavior of the vertical velocity is not at all obvious from any of the equations in the slides.  Namely, how does the vertical component $w$ respond to surface slope and thickness in the non-sliding SIA model?  \emph{Answer:}  It acts diffusively.

In fact, after a calculation similar to that in exercise \textbf{5}, but for the surface value of the vertical component $w$, the whole ice dynamical term in the SKE has the following expression within the non-sliding, non-penetrating planar SIA model:
\begin{align*}
\bu \cdot \bn_s &= \left(u,w\right) \cdot \left(-\frac{\partial s}{\partial x},1\right) = - u \frac{\partial s}{\partial x} + w \\
    &= \frac{1}{2}  A (\rho g)^3 H^{4} \left|\frac{\partial s}{\partial x}\right|^{4} + \frac{\partial}{\partial x} \left(\frac{2}{5}  A (\rho g)^3 H^{5} \left|\frac{\partial s}{\partial x}\right|^{2} \frac{\partial s}{\partial x}\right)
\end{align*}
The $w$ contribution to the SKE is therefore of the form
	$$w = \frac{\partial}{\partial x} \left(D \frac{\partial s}{\partial x}\right)$$
where the diffusivity $D>0$ has a highly nonlinear power-law form.  (\emph{There is a 5th power on thickness and a 2nd power on the slope!})  This diffusive form means, roughly speaking, that a smooth maximum of $s$ tends to be decreased ($\partial s/\partial t < 0$) over time.  Correspondingly, a smooth minimum in $s$ will be increased.

It follows that, in the non-sliding SIA, the velocity can be summarized as \emph{downhill, but with a vertical component that tends to damp-out surface bumps}.  This allows a very informal method of sketching how such velocities will look.

By contrast, it requires more tuition than I have, or have seen, to see at a glance what the surface velocity from the non-sliding Stokes model will look like.  At each location $x$ the velocity has a highly-nontrivial response to the longitudinal stresses transmitted from neighboring ice.  This response is superimposed over a large-scale response to the direction of gravity relative to the general surface slope.  Over several ice thicknesses or more, the Stokes velocity is downhill and damping, just like the SIA model,\footnote{The SIA model is not lying.  Rather, it only makes sense to use the SIA when you are seeking the response to surface and thickness features over a horizontal range of several ice thickness.} but the local velocity response will also depend on longitudinal/membrane effects.  In practice, running a numerical Stokes model for a given geometry would seem to be helpful or even necessary when thinking about the surface dynamical response, and this statement is even more true if there is sliding.


\prob{8}  For the ice thickness $H=s-b$, the inequality equivalent to $s\ge b$ is $H\ge 0$.

Many geophysical flow problems involve a layer of viscous fluid or other material on some substrate of a different character.  The corresponding models describe (parameterize) the layer geometry using thickness or surface elevation.  Examples of these are ocean models, sea ice models, lahar flows, avalanches, and snow fall itself.  Subglacial hydrology offers a rich selection of models, some using water-layer thickness, for example thin films and (spatially-averaged) linked-cavities, while others describe flow geometry in different ways, for instance cross-sectional area $S$ for conduits.  (Note that ``$S\ge 0$'' is relevant to conduits which freeze shut.)

For ocean models the water thickness goes to zero at the ocean shore.  These models can generally ``get away'' with simple wetting/drying mesh schemes because the shoreline location is not highly dynamic or relevant to larger climatic or weather questions.  Exceptions include estuarine and tsunami run-up models.  The loss of the Aral Sea is a climate-relevant exception, in which the ``$H\ge 0$'' constraint is an active participant in determining the model domain.

In summary, glacier, subglacial hydrology, and sea ice models all track layer geometry using a geometry description subject to an inequality constraint, essential a ``$H\ge 0$'' constraint.  In cases where the time scales of surface mass processes and flow are comparable, a margin (terminus, edge, \dots) shape is generated via the balance of an ablation process acting simultaneously with a flow (or elastic displacement, etc.) dynamical process.  The inequality constraint must be actively enforced to correctly determine the modeled shape of the margin.

On the other hand, for practical atmospheric models of the Earth, relevant to weather or climate, the ``$H\ge 0$'' constraint is irrelevant because an atmosphere is always present.  We fortunately never reach an atmospheric margin while traveling on the earth's surface!

\end{document}
